% 第5章 5.3.2：高斯核（sigma=0.15）、观测点 Z={0.1,0.35,0.6,0.9} 的幂函数 P_Z(t)=||K_t-P K_t||_H，
% 以及前一图目标的余项 |r(t)| 与上界 P_Z(t)||r||_H（||r||_H=0.9106）。
% 数据 power.dat 由 ../c5-observation-residual/generate.py 计算；max P_Z=0.5932，max|r|=0.4149。
\begin{tikzpicture}[font=\figurefont\footnotesize]
\begin{axis}[width=82mm,height=46mm,axis lines=left,xmin=0,xmax=1,ymin=0,ymax=0.72,
  xtick={0,0.1,0.35,0.6,0.9,1},xticklabels={$0$,$x_1$,$x_2$,$x_3$,$x_4$,$1$},ytick={0,0.2,0.4,0.6},
  yticklabel style={/pgf/number format/fixed},
  x axis line style={draw=FigureStroke,line width=.6pt},y axis line style={draw=FigureStroke,line width=.6pt},
  tick style={draw=FigureStroke},tick label style={font=\figurefont\scriptsize,text=FigureInk},
  label style={font=\figurefont\scriptsize},xlabel={$t$},xlabel style={at={(axis description cs:1,0)},anchor=north west,yshift=2mm},
  xmajorgrids,ymajorgrids,grid style={draw=FigureGrid,line width=.4pt},
  legend style={draw=none,fill=white,fill opacity=.85,text opacity=1,font=\figurefont\scriptsize,at={(0.5,1.02)},anchor=south,legend columns=3,/tikz/every even column/.append style={column sep=3mm}}]
\addplot[draw=FigureBlue,line width=1.1pt] table[x=t,y=PZ] {figures/02-abstract-mathematics/03-function-spaces/tikz/c5-power-function/power.dat};
\addplot[draw=FigureGreen,line width=1.1pt,dashed] table[x=t,y=bound] {figures/02-abstract-mathematics/03-function-spaces/tikz/c5-power-function/power.dat};
\addplot[draw=FigureYellow,line width=1.1pt,dash dot] table[x=t,y=absr] {figures/02-abstract-mathematics/03-function-spaces/tikz/c5-power-function/power.dat};
\addplot[only marks,mark=*,mark size=1.5pt,FigureInk,forget plot] coordinates {(0.1,0) (0.35,0) (0.6,0) (0.9,0)};
\legend{$P_Z(t)$,$P_Z(t)\|r\|_{H}$,$|r(t)|$}
\end{axis}
\end{tikzpicture}
